\answer{ex:19:1}{(a)~Uma sinapse dá $6.7\mV$ ($R=66.7$~M$\Omega$), então para chegar ao limiar são necessárias no mínimo $4$ (três bastam só assintoticamente); $8$ em $10\ms$; $14$ em $5\ms$. (b)~$0.1\ms\ll\tau_m$; a correção pelo vazamento é $\sim0.25\%$.}
\answer{ex:19:2}{(a)~$1.75\cdot10^{10}$ (um quarto dos misturadores responde a tudo), $4.4\cdot10^9$ e $0.06$ (com $k=40$ quase nunca responde). (b)~$10^3$ vezes: $2\cdot10^5$ contra $200$.}
\answer{ex:19:3}{(a)~$C(2n,n)=2^{2n-1}$ pela simetria dos coeficientes binomiais. (b)~$0.93$ e $0.09$; a largura da transição é $\sim\sqrt{2n}$ em $P$, a relativa, $\sim1/\sqrt n$.}
\answer{ex:19:4}{Com um só dendrito, $V_{\text{corpo}}<\kappa E_E<\theta$ para qualquer número de entradas; com $g_0=0.5g_L$ são necessárias $n\ge3$ em cada dendrito.}
\answer{ex:19:5}{$I\ge C\sigma_V/\sigma_t=15$~nA, isto é, $150$ sinapses síncronas, o que é irrealista; ao neurônio pequeno bastaria $1$~nA, e com a sinapse de $4$~nA o jitter é de $5$~\textmu s.}
\answer{ex:19:6}{Máximo na ponta próxima: $0.03\,\tau_m$ e $0.97$ ($L=0.2$); $0.32\,\tau_m$ e $0.66$ ($L=1$); $0.79\,\tau_m$ e $0.33$ ($L=2$). A estimativa~\eqref{eq:dendriteload} pela capacitância total só vale para $L\ll1$.}
\answer{ex:19:7}{Questão aberta. Com $m$ fixo, o tempo de resposta cresce linearmente com o $P$ memorizado; com fração fixa $m=\varphi n$, ele deixa de depender de $n$, e a lentidão do neurônio grande é consequência da soma de muitas entradas, não do tamanho em si.}
